%% notes-tex/019-simb-multimodal/sections/4-betaxanthin.tex
%% [[experiments.020-cachera-betaxanthin.merzbacher-comparison]]
%% SOURCE: every number in this section is written by
%% experiments/020-cachera-betaxanthin/scripts/plot_merzbacher_comparison.py into
%% experiments/020-cachera-betaxanthin/results/merzbacher_comparison_figures.json, or by
%% experiments/019-simb-multimodal/scripts/pigment_noise_ceiling.py.

\section{Betaxanthin production\secstatus{todo}}
\label{sec:betaxanthin}

\subsection{The task, its ceiling, and the two names that get confused\secstatus{todo}}

\term{Cachera} is a genome-wide betaxanthin deletion screen
(doi:10.1093/nar/gkad656): the four-gene betalain cassette \gene{CYP76AD1}, \gene{DOD},
\gene{ARO4}$^{\mathrm{K229L}}$ and \gene{ARO7}$^{\mathrm{G141S}}$ was transferred by CRI-SPA
into each strain of the yeast knockout collection and per-colony fluorescence read as a
production proxy. It supplies 4,735 single deletions and is the ground truth here.
\term{Flux Cone Learning} (FCL) is Merzbacher et al.'s method, which learns from flux
samples drawn through a genome-scale metabolic model;
\file{RandomForestClassifier_Resampled} is its best variant. The comparison is CGT
against FCL, both scored on the Cachera screen.

Betaxanthin is the only phenotype in this document with a high ceiling. Its records carry a
per-strain standard error over a median of 15 colonies, giving a reliability of $0.836$ and
a Pearson ceiling of $\mathbf{0.914}$.
\src{experiments/019-simb-multimodal/results/pigment_noise_ceiling.json}

\tcfigfit[0.62]{figures/betaxanthin_noise_ceiling.pdf}{Reliability of the betaxanthin
target. Unlike every other phenotype here, Cachera releases a per-strain standard error over
a median of 15 colonies, so the noise term is measured directly rather than inferred from
replicate structure. That is what puts this ceiling at $0.914$ and makes betaxanthin the one
strand where a regression target is clearly
supportable.}{fig:bx-ceiling}

\subsection{Three betaxanthin numbers, two scoring rules, and one honest comparison\secstatus{todo}}
\label{sec:bx-two-numbers}

\begin{table}[htbp]
\centering\small
%% SOURCE: results/analysis_002_noise_and_lambda.json; results/round_leaderboards.csv;
%% OBJECTIVE_METRIC in experiments/020-cachera-betaxanthin/scripts/optuna_metabolism_sweep.py
\begin{tabular}{llrrl}
\toprule
number & where it comes from & scoring rule & $n_{\mathrm{train}}$ & split \\
\midrule
$\mathbf{0.4301}$ & Optuna study \file{betaxanthin_002} & raw per-epoch max & 4{,}235 & ordinary ratio split \\
$\mathbf{0.4015}$ & the same runs, in W\&B & \texttt{roll\_max} & 4{,}235 & ordinary ratio split \\
$\mathbf{0.3705}$ & \file{torchcell_020_betaxanthin_v4} & \texttt{roll\_max} & 3{,}698 & Merzbacher test genes pinned out \\
\bottomrule
\end{tabular}
\caption[]{The first two rows are the \textbf{same runs under two scoring rules}, which is
what made the previous revision confusing. The Optuna objective is
\file{val/betaxanthin/pearson_per_feature_max}, a running maximum over epochs; the
leaderboard's \texttt{roll\_max} first averages five neighboring epochs, and smoothing can
only lower a maximum. Only rows two and three are comparable, and their difference,
$0.4015$ against $0.3705$, is what holding out a competitor's test set costs. Against the
$0.914$ ceiling the best is $47\,\%$ realized on the raw rule and $44\,\%$ on the smoothed
one.}
\label{tab:bx-two-numbers}
\end{table}

\notesec{On whether $0.4301$ is a plateau or one lucky trial, since the previous phrasing of
this was unreadable} A \term{trial} here is one Optuna trial, meaning one hyperparameter
configuration trained once. It is not a data split: the split is pinned across the whole
study, so trials differ by hyperparameters and by weight initialization only. Study
\file{betaxanthin_002} ran 37 of them. Two things follow, and they point in opposite
directions.

\begin{itemize}
\item \textbf{The top is flat, so the winner is not a spike.} The best five trials read
  $0.4301$, $0.4234$, $0.4216$, $0.4211$ and $0.4095$, a spread of $0.0206$. Re-running an
  \emph{identical} configuration moves the objective by $\sigma = 0.0302$
  (\cref{tab:replicate-noise}), so those five sit inside two thirds of one standard
  deviation of each other. They are statistically indistinguishable, which means the winner
  is one of five equivalent configurations rather than a single fortunate run.
\item \textbf{But taking the maximum of 37 noisy trials is still biased upward.} That is
  \term{selection inflation}, estimated at $0.064$ here. So a configuration chosen in
  advance, without seeing the study, should be expected to score around $0.366$ rather than
  $0.430$.
\end{itemize}

Both are true at once: the plateau says the number is not a fluke of one run, and the
inflation says it is not what a fresh configuration would give.

\notesec{The comparison number is lower than the best number on purpose} The head-to-head
below runs on \file{020_v4}, whose split \textbf{pins the 640 Merzbacher test genes out of
training}. That costs 537 training strains, and it costs score. Using the higher-scoring
model instead would mean scoring on genes it had trained on, which is the one thing that
would void the comparison. The cell carried forward is \file{s09_L6_maskon_lr0.0001_yj},
selected by validation at $0.3639$ and never by its score on the comparison genes.

\notesec{The 537 strains are not recoverable, and the fix is a reporting change that needs
no retraining} The natural proposal is to train on everything outside the comparison list,
carve validation and test from that same remainder, and evaluate twice. Sized against the
data, that proposal is numerically the split \file{020_v4} already ran: it gives
$3{,}682/299/295$ plus the 639 comparison genes, against \file{020_v4}'s actual
$3{,}698/286/931$, where the 931 is 639 comparison genes plus 292 of our own. Holding the
comparison genes out of training is what costs the 537, under any scheme, and no
rearrangement returns them.

\begin{table}[htbp]
\centering\small
%% SOURCE: experiments/020-cachera-betaxanthin/scripts/fcl_resplit_sizing_check.py
\begin{tabular}{lrrr}
\toprule
arm & train & validation & test \\
\midrule
unpinned (\file{torchcell_020_betaxanthin}, \file{_v3}) & 4{,}235 & 340 & 340 \\
pinned (\file{torchcell_020_betaxanthin_v4}) & 3{,}698 & 286 & 931 \\
the proposed re-split, sized & 3{,}682 & 299 & 295 + 639 \\
\bottomrule
\end{tabular}
\caption[]{Sizing the proposed re-split against what was run. The bookkeeping on the pinned
arm closes exactly: train $-537$, validation $-54$, test $+591 = 537 + 54$, the other 48
pinned genes having already been in the unpinned test set. The realized split ratio is
$86/7/7$ rather than the nominal $80/10/10$, because assignment is per-key stratified; the
config comment estimating ``roughly 511'' lost genes used the nominal ratio and understates
the true 537 by 26.}
\label{tab:bx-resplit}
\end{table}

\textbf{What the proposal does buy, and it is worth taking, is the second evaluation.} The
existing 931-record test set is already 639 comparison genes plus 292 of our own. Splitting
it and reporting both numbers separates ``how we do against FCL on their genes'' from ``how
the model does on a clean held-out sample'', and it runs off the existing \file{020_v4}
checkpoints with no retraining at all. Enlarging our own test set beyond those 292 is the
only part that costs anything, at roughly a further 137 training strains for a 429-gene own
test set.

\notesec{One name-resolution defect, now triangulated and resolved} Of the 640 comparison
genes, 630 carry a betaxanthin measurement, and the pin file holds 639. The screen reports
one gene as \gene{PPA1}, an ambiguous alias naming both \gene{IPP1} (\gene{YBR011C}), which
is essential, and \gene{VMA16} (\gene{YHR026W}), which is not. Neither owns \gene{PPA1} as
its standard name, so the resolver's standard-name-before-alias rule cannot break the tie and
correctly returns \texttt{AMBIGUOUS}. The raw screen does not settle it either: its 27
columns carry a gene name and colony statistics, with no systematic ORF, plate, position or
barcode.

Five independent lines were run, and they agree.
\src{experiments/020-cachera-betaxanthin/scripts/triangulate\_ppa1\_alias.py}

\begin{itemize}
\item \textbf{The strain is alive, which is decisive on its own.} The \gene{PPA1} row records
  \textbf{15 colonies at 24 h and 15 at 48 h}, with mean areas of $106.5$ and $118.5$. A
  haploid \gene{ipp1}$\Delta$ is inviable and yields no colonies, so no row for it can exist.
  This is the strain's own growth data and does not depend on any claim about which
  collection was screened. The colony is small, in the $1.6$th percentile of screen area,
  which is what a severe but survivable defect looks like.
\item \textbf{The screen contains essentially no essential genes}, on $1{,}140$
  counterexamples rather than the six previously cited. It covers $4{,}704$ of $5{,}467$
  non-essential genes ($86\,\%$) and $29$ of $1{,}140$ SGD inviable-null genes ($2.5\,\%$),
  and those 29 are autophagy genes and paralogs whose inviability is conditional, all of
  which grew.
\item \textbf{The V-ATPase family has a \gene{VMA16}-shaped hole.} Thirteen of the fifteen
  \gene{VMA} genes are in the screen; the only absentees are \gene{VMA9} and \gene{VMA16}.
  The file names genes by alias constantly (\gene{VMA1} as \gene{TFP1}, \gene{VMA3} as
  \gene{CUP5}), so an alias row for \gene{VMA16} is exactly what one expects.
\item \textbf{Its value sits on top of the V-ATPase cluster.} Betaxanthin $-1.566$ against
  \gene{VMA13} $-2.19$, \gene{VMA3} $-1.55$, \gene{VMA22} $-1.54$, \gene{VMA6} $-1.44$, on a
  screen median of $-0.006$. Soft corroboration, carrying nothing alone.
\item \textbf{Neighbor evidence is neutral here}, unlike the \gene{FEN1} case: neither
  candidate appears elsewhere in the screen under another name, so the ``one row per gene''
  argument that decides \gene{FEN1} does not apply.
\end{itemize}

\textbf{The \gene{PPA1} row is \gene{YHR026W} (\gene{VMA16}), and the split builder is
right.} Nothing points to \gene{IPP1} except the build itself.

\notesec{The training build has it wrong, and the head-to-head does not} The defect is real
but narrower than feared. In both the loader LMDB and the \file{fig6_pigment_transfer}
training build, \gene{YBR011C} carries the betaxanthin value $-1.5664$ and \gene{YHR026W}
carries none. The cause is verified in code rather than inferred: the build predates the
layered resolver, and the loader then took \file{candidates[0]} from the alias map, which
returns \gene{YBR011C} for \gene{PPA1}. The same line got \gene{FEN1} right by luck.

\textbf{The comparison is not corrupted}, because \file{build_merzbacher_split.py} applies
the screen-local disambiguation when it reads the raw file, so the comparison's ground truth
for \gene{YHR026W} is the \gene{PPA1} value and \gene{YBR011C} is the gene dropped from 640
to 639. What remains is one false training label out of 4,735: the model was trained on a
betaxanthin value attributed to a gene whose deletion strain cannot exist, and scored on the
same measurement under its correct ORF.

\textbf{A plain rebuild does not fix this, it loses the value.} The current loader rejects
\texttt{AMBIGUOUS}, so it would drop the row and leave neither ORF labeled. The
disambiguation table lives in the split builder and needs to reach the loader, which is
issue \#195 territory and lands with the next full graph build.

\notesec{And the same alias broke Merzbacher's pipeline too} Their released 640-gene test
list contains \emph{both} \gene{YBR011C} and \gene{YHR026W}, each labeled low, so their code
expanded the alias to both candidates rather than disambiguating. It is an independent
instance of the same failure, and it is worth knowing that 5 of their 640 test genes are
SGD inviable-null in what is otherwise a non-essential list.

\subsection{The published accuracy comparison is empty for both methods\secstatus{todo}}

FCL reports gene-level accuracy $0.700$ against a majority-class rate of $0.673$, achieved
by calling $607$ of $640$ genes medium, which is $94.8\,\%$. Our absolute binning sits at
$0.681$ doing the same thing. Forced to the true class marginal, FCL drops to $0.556$ and
CGT to $0.549$.

\notesec{The structural point} Every model that finds more high producers has worse overall
accuracy, because the only way to call more highs is to stop calling everything medium.
Accuracy selects against the capability the task is about, so it should not be the primary
metric in any writeup, including ours.

\subsection{Both methods have real signal, and it lives in the tail\secstatus{todo}}

Of the top $k$ genes each method nominates, the fraction that are truly high producers,
against a base rate of $0.156$:

\begin{table}[htbp]
\centering\small
%% SOURCE: results/merzbacher_comparison_figures.json, fig3_precision_at_k
\begin{tabular}{lccc}
\toprule
model & $k = 10$ & $k = 25$ & $k = 50$ \\
\midrule
FCL RF & 8 ($5.1\times$, $p = 10^{-5}$) & 19 ($4.9\times$, $p = 8\times10^{-12}$) & 20 ($2.6\times$, $p = 10^{-5}$) \\
CGT    & \textbf{9} ($5.8\times$, $p = 4\times10^{-7}$) & 18 ($4.6\times$, $p = 10^{-10}$) & \textbf{24} ($3.1\times$, $p = 2\times10^{-8}$) \\
\bottomrule
\end{tabular}
\caption[]{Precision at $k$ for high producers. Nine of CGT's top ten are genuine high
producers. Both curves approach the base rate by $k = 150$.}
\label{tab:precision-at-k}
\end{table}

Median percentile rank by true class is flat for both methods, at 54/49/56 for FCL and
48/51/52 for CGT across low, medium and high. Read with the table above, that says the
signal is real and confined to the tail, not that there is none: neither method separates a
typical high producer from a typical low one, and both are far better than chance at the very
top of their own ranking.

\tcfigfit{figures/merzbacher_fig3_precision_at_k_2026-08-02-23-43-08.pdf}{Precision at $k$
for high producers. (a) the comparison: FCL against the CGT cell chosen by validation, with
the nine unselected grid cells drawn as a shaded envelope rather than as nine competing
lines. The base rate of $0.156$ is drawn and labeled, and it is what a random pick of $k$
genes would return. Nine of CGT's top ten are genuine high producers, and both methods reach
the base rate by $k \approx 150$, which is where the tail signal is exhausted. (b) the same
ten cells drawn individually, one color each, which is what answers whether some cells sit
below the base rate early: \textbf{three of the ten do at $k = 50$}. Read crossings at very
small $k$ with care, since at $k = 10$ precision moves in steps of $0.1$ and a crossing is a
single gene. The spread across cells is far larger than the gap between the two methods,
which is the hyperparameter sensitivity of \cref{sec:terms} rather than a split
effect.}{fig:precision-at-k}

\tcfigfit{figures/merzbacher_fig1_scatter_spread_2026-08-02-23-43-08.pdf}{Per-gene spread
over the 639 shared genes. Both models emit a nearly constant value: $80\,\%$ of FCL's
genes fall inside an expected-class band of $0.22$ on a $0$ to $2$ scale, $80\,\%$ of CGT's
inside $0.33$ $z$-units. Panel c shows the disagreement concretely, with each method
recovering 29 of the 100 true high producers and sharing 9.}{fig:bx-spread}

\subsection{The two methods find different genes\secstatus{todo}}

Rank-matched into the same 100 high slots, FCL recovers 29 of the 100 true high producers
and CGT recovers 29, and they share 9. The union reaches 49, so running both finds $69\,\%$
more high producers than running either.

\notesec{The two methods' predictions are anti-correlated, and here is what that sentence
means} Compare their predicted rankings over the 639 shared genes directly, ignoring the
truth. Two methods that are each independently right about the same target should agree
somewhat, and the amount is calculable: if one correlates with the truth at $+0.105$ and the
other at $+0.039$, then agreeing only through the truth would put their mutual correlation
at about $0.105 \times 0.039 = +0.004$, essentially zero. The observed correlation between
them is $\mathbf{-0.128}$, which at $n = 639$ is $z = -3.2$: they disagree substantially
more than two independent methods would.

\notesec{No mechanism is claimed for the negative sign} It could be a genuine
anti-relationship between what the two feature spaces encode, or an artifact of how each was
fit; nothing here distinguishes those. What \emph{is} claimed is the operational
consequence, which does not depend on the sign's explanation: because the two find largely
different genes, their recall adds rather than overlaps, so running both is worth more than
picking the better one. That fact is invisible in every aggregate statistic, all of which
report the two as tied. The 51 high producers that neither finds is the honest ceiling on
this comparison.

\subsection{The real difference is coverage\secstatus{todo}}

yeast-GEM contains 1,161 genes, of which 915 appear in the Cachera screen. That is
$19.4\,\%$ of its 4,721 deletions. A flux-based method has features for those and for
nothing else.

\begin{table}[htbp]
\centering\small
%% SOURCE: results/merzbacher_comparison_figures.json, fig8_screen_coverage
\begin{tabular}{lrr}
\toprule
population & in yeast-GEM & outside \\
\midrule
every deletion in the screen & 915 (19\,\%) & \textbf{3{,}806 (81\,\%)} \\
top 100 measured producers & 32 (32\,\%) & \textbf{68 (68\,\%)} \\
all high producers & 60 (27\,\%) & \textbf{163 (73\,\%)} \\
\bottomrule
\end{tabular}
\caption[]{Metabolic genes are enriched among high producers, at $26.9\,\%$ of the highs
against $19.4\,\%$ of the screen, a $1.4\times$ enrichment that is exactly what the biology
predicts. It is not enough to matter: restricting to metabolism buys that concentration and
costs three quarters of the hits.}
\label{tab:gem-coverage}
\end{table}

\tcfigfit{figures/merzbacher_fig8_screen_coverage_2026-08-02-23-43-08.pdf}{Measured data
only, no model and no prediction: how little of the betaxanthin screen a flux-based method
can reach.}{fig:gem-coverage}

\notesec{How the class labels are derived, since every metric below depends on it} The
Cachera screen releases a continuous production value, and FCL's task is a three-class
one, so a rule is needed to turn measurements into low, medium and high. FCL's rule:
min-max scale production onto $[0, 1]$ over their gene set and cut at $0.40$ and $0.65$.
\textbf{We apply their thresholds rather than inventing our own}, which is what makes the
two panels comparable. Two consequences follow and both are stated rather than hidden.

\begin{itemize}
\item \textbf{Applying their rule to our larger copy of the same screen does not reproduce
  their class counts.} It gives $107/476/56$ against their $109/431/100$, and $18.8\,\%$ of
  genes land in a different bin than our measured values would give. A model predicting our
  measurement \emph{perfectly} is therefore marked wrong on about a fifth of genes, so their
  labels are a headroom bound on any accuracy-style metric, ours included.
\item \textbf{They never labeled the non-GEM genes at all}, because those genes are outside
  their method's reach. We derive labels there with the identical rule, which is the only
  way that panel exists. The comparison to make is therefore panel against panel rather than
  number against number, and \cref{fig:capability-gap} is drawn that way.
\end{itemize}

\notesec{Why every metric here is rank-based or re-binned} We produce a regression and they
produce a three-class classifier, so nothing can be compared directly. Two repairs are used
throughout, and the choice between them is not arbitrary. \textbf{Rank-based metrics}
(Spearman, precision at $k$, AUC) need no thresholds at all and so cannot be affected by the
binning disagreement above; they are preferred wherever available. Where a class metric is
unavoidable, our continuous predictions are binned \textbf{with thresholds fit on the
training split only}, never on the test genes, so the binning cannot borrow information from
the evaluation. Neither repair fixes the $18.8\,\%$ label disagreement, which sits in the
ground truth rather than in either model.

On held-out test genes split by yeast-GEM membership, with labels derived the same way on
both sides, CGT's top 25 recovers 17 of 48 true highs inside the model ($9.3\times$ over
base rate, $p = 2\times10^{-15}$) and 9 of 21 outside it ($4.4\times$,
$p = 2\times10^{-5}$).

\tcfigfit{figures/merzbacher_fig9_metabolic_vs_nonmetabolic_2026-08-02-23-43-08.pdf}{CGT
predicts high betaxanthin producers among deletions a flux-based method cannot represent.
Held-out test genes split by yeast-GEM membership, identical derived labels on both sides,
circled points are CGT's top 25. Flux Cone Learning has no prediction for the right
panel.}{fig:capability-gap}
FCL's prediction for the outside set is not poor, it is undefined: flux sampling yields no
feature for a deletion the metabolic model does not contain. The FCL test set is $639$ of
$640$ inside yeast-GEM; of our other test genes only $21$ of $294$ are.

\notesec{One qualification that must travel with this claim} ``Inside'' and ``outside'' here
mean inside or outside \textbf{yeast-GEM}, the genome-scale metabolic model, not inside or
outside anyone's training set. CGT's rank correlation with measured production is
\emph{higher} on the genes yeast-GEM does not contain, $+0.270$ against $+0.124$. That is the
opposite of what one would expect if metabolic genes were the easy ones, so it is worth
testing rather than asserting: the two correlations are measured on different,
non-overlapping gene sets, and comparing two correlations requires first mapping each
through Fisher's transform, which turns a correlation into a quantity whose sampling
distribution is approximately normal with a known variance. Doing that gives $z = 2.07$,
$p = 0.039$: a real difference, but marginal, and on populations that differ in more than
their yeast-GEM membership. \textbf{The safe reading is that CGT is at least as good outside
the metabolic model as inside it}, which is all the coverage argument needs.

\subsection{Where the comparison is soft\secstatus{todo}}

\begin{itemize}
\item \textbf{The selection is asymmetric, and it favors FCL. This is now verified rather
  than suspected, and it is stronger than the previous wording.} Revision 2 said their
  validation folds ``appear to be drawn from the same 640 genes as their test list''. From
  their Zenodo release: \file{yeast_production_validation_split.csv} assigns a fold to each
  of \textbf{exactly the same 640 systematic ORFs} that
  \file{yeast_production_test_split.csv} names. The intersection is 640 and both set
  differences are empty, in five folds of 129/128/128/128/127. Their training loop then
  computes \file{train_names = set(nontest_names) - set(val_names)}, which is a no-op
  because those genes were never in the non-test pool, so every fold trained on the full
  complement and validated on a 128-gene slice of the test set. Model selection across their
  16 variants therefore ran on test data.

  Their own CHO essentiality release does \emph{not} do this: there the validation genes are
  exactly the 1,833 flagged as non-test, sharing nothing with the 457 test genes. So the
  yeast leak is a deviation from their practice inside the same release rather than a naming
  convention, which is what makes it worth stating as a fact. And the paper alone cannot
  settle it: its Methods never mention a validation set and the test count is internally
  inconsistent (659 in the main text, 649 in Table S6, 640 in the release), so only the
  released files decide it.

  Ours, by contrast, is the validation-selected cell with the comparison genes pinned out of
  training. \textbf{So where the two land level, we are level with an opponent selected on
  the test set.}
\item \textbf{Validation Pearson does not track test AUC, and this is the hyperparameter
  sensitivity again rather than a split effect.} Four of ten CGT cells beat FCL's AUC of
  $0.570$ and the best reaches $0.695$, while the validation-selected cell reads $0.557$.
  The runner-up on validation, $0.3528$ against $0.3639$, has an AUC $0.14$ higher. Every
  one of those cells sees the same split; what differs is the hyperparameters
  (\cref{sec:terms}). So the honest statement is that \textbf{our own model-selection
  criterion does not rank cells the way the comparison metric does}, which makes neither
  ``CGT's AUC is 0.695'' nor ``CGT's AUC is 0.557'' quotable on its own. Selecting on
  validation and reporting whatever AUC follows is the defensible protocol, and it is what
  was done; the spread is reported beside it rather than being selected from.
\item \textbf{Their labels are a headroom bound}, as detailed above: $18.8\,\%$ of genes bin
  differently under their rule than our measured values would give, so a model predicting
  our measurement perfectly is marked wrong on about a fifth of genes.
\item \textbf{Regression against classification.} They tried regression and abandoned it,
  in their words ``challenging with the limited number of knockouts at the high and low
  ends.'' Our ceiling of $0.914$ makes regression available to us. The comparison is
  therefore between a regressor and a three-class classifier, which is why every metric
  above is either rank-based or computed after binning our predictions with thresholds fit
  on the training split only.
\item \textbf{The build is stale} with respect to the gene-name resolver (issue \#195),
  which the split builder avoids by working from the raw screen, and which the training data
  still inherits until a rebuild.
\end{itemize}

\begin{keybox}
\textbf{The betaxanthin case is about coverage, and it does not need to beat anyone on
accuracy.} Where both methods can be scored, CGT and Flux Cone Learning are tied, both with
tail-confined signal at roughly $5\times$ enrichment in the top 25, and they find largely
different genes, so their recall adds: 29 each of 100 true high producers, sharing 9,
reaching 49 together. The asymmetry favors them, since their model was selected on its own
test set and ours on validation. \textbf{The argument that carries is the other one}: a
flux-based method has features for $19\,\%$ of the screen and misses $73\,\%$ of its high
producers, while CGT predicts high producers in the $81\,\%$ it cannot represent at
$4.4\times$ over base rate, and scores at least as well there as inside the metabolic model.
\end{keybox}

\notesec{Cannot say} That CGT beats FCL on any single aggregate. That the rank
anti-correlation means anything mechanistically. That the absolute top-100 membership is a
deployable call: the rank-matched comparison works by handing \emph{both} methods the true
number of high producers and asking which genes each puts in those 100 slots. That is the
right way to compare two methods, because neither can win by predicting more or fewer highs
than exist. But it is not available in deployment, where nobody knows how many high
producers a screen contains, so the top-100 precision measured here is an upper bound on
what the same model would deliver prospectively.
